\documentclass{article}
\usepackage{amsfonts}
\usepackage{graphicx} % Required for inserting images

\title{Taller 6}
\author{Gonzalo Balladares}
\date{12 de Marzo de 2025}

\begin{document}

\maketitle

\section{Pregunta 1}

\begin{enumerate}
    \item \textbf{Serial} S no es serial, ya que considerando las transacciones no se ejecutan de forma aislada, se comienza a ejecutar la T2 sin haber terminado la T1, con lo que la schedule no es serial.
    \item \textbf{conflict-serializable} No es conflict serializable por el W1(a) y el W2(a).
    \item \textbf{serializable} No se puede determinar al no saber cuáles son las operaciones realizadas en W(a)
\end{enumerate}

\section{Pregunta 2}
Si es conflict serializable, a continuación se indican las permutaciones para serializar la schedule:
\begin{enumerate}
    \item (R1(X), W1(X)), (W3(Y), W3(Z))
    \item R2(Z), (R1(Y), W1(Y))
\end{enumerate}

Con estas permutaciones se ejecuta T3 en su totalidad antes de ejecutarse T1, quedando T3 hacia el final, sin hbaer ningún conflicto.
Nótese que T3 corresponde a la columna derecha de la tabla del enunciado, que aparece como T2 junto con la columna central.

\section{Pregunta 3}
El schedule no es conflict-serializable. Se podría hacer la permutación de [(R1(X), W1(X)), (W3(Y), W3(Z))], con lo que las transacciones T1 y T3 se ejecutan de forma aislada, pero últimas 2 operaciones -pertenecientes a T2-, se permutar de forma tal que T2 también se ejecute de forma aislada, pues (R1(X), W1(X)) tiene conflicto con (R2(X), W2(X)), con lo que la schedule no es conflict-serializable.

\section{Ejemplo}

\section{Ejemplo de Schedule Conflict-Serializable}

\begin{table}[h!]
\centering
\begin{tabular}{|c|c|c|}
\hline
\textbf{T1}       & \textbf{T2}       & \textbf{T3}       \\ \hline
R1(X)             &                   &                   \\ \hline
W1(X)             & R2(X)             &                   \\ \hline
                  & W2(Y)             &                   \\ \hline
                  &                   & R3(Y)             \\ \hline
                  &                   & W3(Z)             \\ \hline
R1(Z)             &                   &                   \\ \hline
W1(Z)             &                   &                   \\ \hline
\end{tabular}
\caption{Schedule con transacciones T1, T2 y T3.}
\label{tab:conflict_serializable_example}
\end{table}

\section{Ejemplo de Schedule Serializable pero No Conflict-Serializable}

\begin{table}[h!]
\centering
\begin{tabular}{|c|c|c|}
\hline
\textbf{T1}       & \textbf{T2}       & \textbf{T3}       \\ \hline
R1(X)             &                   &                   \\ \hline
                  & R2(X)             &                   \\ \hline
W1(X)             &                   &                   \\ \hline
                  &                   & W3(X)             \\ \hline
                  & W2(Y)             &                   \\ \hline
                  &                   & R3(Y)             \\ \hline
\end{tabular}
\caption{Schedule con transacciones T1, T2 y T3.}
\label{tab:serializable_not_conflict_serializable}
\end{table}

\subsection{Conflictos en el Schedule}
En este schedule, no se puede determinar un orden serial equivalente utilizando únicamente conflictos, ya que:
\begin{itemize}
    \item \textbf{Conflicto 1:} \texttt{W1(X)} (T1) y \texttt{R2(X)} (T2) tienen un conflicto de \textbf{lectura-escritura}.
    \item \textbf{Conflicto 2:} \texttt{W3(X)} (T3) y \texttt{W1(X)} (T1) tienen un conflicto de \textbf{escritura-escritura}.
    \item \textbf{Conflicto 3:} \texttt{W2(Y)} (T2) y \texttt{R3(Y)} (T3) tienen un conflicto de \textbf{escritura-lectura}.
\end{itemize}

\subsection{Por qué no es Conflict-Serializable}
El grafo de precedencia (precedence graph) para este schedule tiene un ciclo:
\[
T1 \to T2 \to T3 \to T1
\]
Esto indica que no es conflict-serializable, ya que no se puede reordenar el schedule para eliminar los conflictos.

\subsection{Por qué es Serializable}
Aunque no es conflict-serializable, este schedule es \textbf{serializable} porque produce el mismo resultado que el siguiente orden serial:
\[
T1 \to T3 \to T2
\]

\subsection{Explicación de la Serialización}
\begin{itemize}
    \item Primero, T1 ejecuta sus operaciones (\texttt{R1(X)}, \texttt{W1(X)}).
    \item Luego, T3 ejecuta sus operaciones (\texttt{W3(X)}, \texttt{R3(Y)}).
    \item Finalmente, T2 ejecuta sus operaciones (\texttt{R2(X)}, \texttt{W2(Y)}).
\end{itemize}

Este orden serial produce el mismo resultado final que el schedule original, por lo que es serializable.
\end{document}
